\documentclass[11pt]{article}
\usepackage{amsmath,amssymb,amsthm,hyperref,geometry}\geometry{margin=1in}
\newtheorem{theorem}{Theorem}\newtheorem{lemma}{Lemma}
\title{Diagram lifting of a presented native interaction}
\date{15 November 2026}\begin{document}\maketitle
\paragraph{Status.} Scoped closure theorem with written proofs. This lifts an \emph{existing} square-functorial presented context/content interaction to a diagram pure base. It is not a new existence criterion for an arbitrary pair of doctrines. The construction is categorical; effects and couplings are optional realizations. It neither reverses native arrows nor commutes a diagram functor with an unproved pseudocoend.

\section{Precise presentation input}
Let $\mathcal C$ be the pure category, with nested universe choices. A context presentation is a fibration $p_K:\mathcal K\to\mathcal C$ with extension $e:\mathcal K\to\mathcal C$ and its valid contextad unit, dependent accumulation, and coherent native/pure-square actions. A content presentation is an opfibration $p_L:\mathcal L\to\mathcal C$ with extension $d:\mathcal L\to\mathcal C$ and its valid contentad operations. Cleavages can be used for notation; their coherence is retained, not declared identity.

Use the original one-sided modules
\[
 M(A,B)=(e_A\downarrow B),\quad N(A,B)=(A\downarrow d_B),\quad H(A,B)=(e_A\downarrow d_B).
\]
An endpoint square over $u:A\to A',v:B\to B'$ is an actual total context arrow $\alpha:P\to P'$ over $u$ and total content arrow $\beta:Q\to Q'$ over $v$, satisfying $d\beta f=f'e\alpha$. In cleavage notation these are $P\to u^*P'$ and $v_!Q\to Q'$, with the corresponding extension arrows inserted.

In addition to an existing full law $\tau:N\odot M\Rightarrow M\odot N$, require its given presentation to have the following explicit form, rather than inferring it from fixed-endpoint components alone. The crossing sends each source representative
\[
 (A,Y,B;\ Q\in\mathcal L_Y,\ a:A\to dQ;\ P\in\mathcal K_Y,\ b:eP\to B)
\tag{R}
\]
to a literal mixed object of $H(A,B)$. It is a functor on the total category of such representatives and their simultaneous endpoint/middle/native squares. Its pure-middle balance comparisons are coherent natural isomorphisms. The four pseudodistributive comparisons and their diagrams hold on these generators and are natural for those total squares. Call this the \emph{square-functorial representative presentation}. It is stronger data than merely knowing a multiplication on isomorphism classes; it is exactly the presentation property to check in an application.

In the pinned original O1 source, the endpoint equation, retained normalization lemma, native-map construction, arbitrary endpoint maps and pure-middle Beck--Chevalley balance are the relevant checks. This theorem reuses that full presentation, not just the assertion that some law exists. It does not rederive O1's independent-input criterion or relax its content hypotheses.

\section{Diagram contexts and content}
Fix a small category $J$. Let the new pure category be $\mathcal C^J$. Set
\[
 p_K^J:\mathcal K^J\to\mathcal C^J,\quad e^J:\mathcal K^J\to\mathcal C^J,
 \qquad p_L^J:\mathcal L^J\to\mathcal C^J,\quad d^J:\mathcal L^J\to\mathcal C^J.
\]
A context over $A:J\to\mathcal C$ is an actual functor $P:J\to\mathcal K$ with $p_KP=A$. Thus each arrow $j\to k$ carries a potentially noninvertible total context map above $A_j\to A_k$, not a cartesian-only or invertible section. Content is similarly an actual diagram in $\mathcal L$. A mixed body is a natural transformation $f:eP\Rightarrow dQ$.

\begin{lemma}[One-sided diagram presentations]
The displayed context projection is a fibration and the content projection an opfibration. Their extensions and one-sided operations give diagram contextad/contentad presentations with all native transformations retained.
\end{lemma}
\begin{proof}
For a natural transformation $u:A\Rightarrow A'$ and a context diagram $P'$ above $A'$, choose cartesian lifts over each $u_j$. The unique factorization property defines the structure arrows of the lifted diagram; naturality of $u$ makes them well-typed. Uniqueness gives identities, composition, and naturality of the cartesian transformation. Componentwise universality proves its universality in the functor category. The opcartesian argument is dual. Unit and dependent accumulation are the given total operations applied componentwise. Their naturality supplies the structure arrows of the result diagram, and their coherence diagrams hold at every component, hence as diagrams of natural transformations. In particular dependent continuation contexts are over the full diagram $eP$, not chosen separately with incompatible arrow actions.
\end{proof}

Write $M^J,N^J,H^J$ for the literal comma modules of these diagram presentations, not for an unproved pointwise exponentiation of the old relative tensors.

\begin{lemma}[Diagram retained normalization]
There is a native module equivalence $M^J\odot_{\mathcal C^J}N^J\simeq H^J$.
\end{lemma}
\begin{proof}
A representative with middle diagram $Y$ and natural bodies $g:eP\Rightarrow Y$, $h:Y\Rightarrow dQ$ maps to $(P,Q,hg)$. Conversely $(P,Q,f)$ has middle diagram $eP$ and representative $(P,\mathrm{id});(Q,f)$. For a native transformation $(\alpha,\beta)$ with $d\beta f=f'e\alpha$, apply the $N^J$ arrow $\beta$, balance along the \emph{whole natural transformation} $e\alpha:eP\Rightarrow eP'$, then apply the $M^J$ arrow $\alpha$. This is the explicit inverse arrow action. Native composition and naturality are the defining balance relations. The composite to $H^J$ is identity; the other composite is naturally isomorphic to identity by balancing along $g$. The triangles are balance identity/composition. Pure endpoint actions commute by the diagram cleavage comparisons. No claim that $(-)^J$ preserves pseudocoends is used.
\end{proof}

\section{The lifted crossing and coherence}
\begin{theorem}[Presented diagram interaction]
A square-functorial representative presentation as above induces a full presented law
\[
 \tau^J:N^J\odot_{\mathcal C^J}M^J\Longrightarrow M^J\odot_{\mathcal C^J}N^J
\]
and the induced relative multiplication on the literal diagram comma family $H^J$. All native diagram arrows, including noninvertible ones, are retained.
\end{theorem}
\begin{proof}
A source representative with middle diagram $Y$ gives a functor from $J$ into the total representative category (R): its arrows comprise the natural input/output bodies, the middle arrows of $Y$, and the actual context/content diagram arrows. Apply the stipulated representative crossing functor. Its output is a functor into the total mixed-square category; equivalently it is context/content diagrams $\widetilde P,\widetilde Q$ with a natural mixed body. Functoriality of that total crossing is what proves the output section identities and composition, not an unjustified pointwise choice of output representatives.

Transformations of source representatives map to transformations of these output diagrams by the same functor. A pure-middle balance along a natural map $Y\Rightarrow Y'$ gives an invertible transformation of output diagrams: its components are the original balance maps, and their naturality for total squares is exactly naturality along every arrow of $J$. Identity/composition relations descend, so this defines a functor to $H^J$ on the relative pseudocoend generators. Compose with the explicit normalization inverse to obtain the displayed target module.

The lifted context/content multiplication and unit comparisons are the original four comparisons at each component. Their total-square naturality makes them transformations between the constructed diagram functors. They are invertible because their components are. Every pseudodistributive coherence equation holds after evaluation at every object of $J$, using the corresponding original equation. Natural transformations in a functor category are equal if all components are equal; therefore the diagram equations hold. Where a tensor representative is involved, the preceding generator-level construction and balance relations establish well-definedness before this pointwise check. Standard relative pseudomonad composition then gives the claimed multiplication on $H^J$ through diagram normalization.
\end{proof}

\section{Verification for the pinned O1 presentation}
The closure theorem does not claim that $\mathcal C^J$ is itself Set or that it independently satisfies the original polynomial display criterion. It lifts the categorical package already constructed by that criterion. Here is the essential total-arrow check, beyond fixed-endpoint naturality.

For a generic O1 crossing write $P\in\mathcal K_B$, $Q\in\mathcal L_B$, $\pi:E=eP\to B$, and output content $\pi^*Q$. Suppose actual arrows $P\to P'$ and $Q\to Q'$ lie above $b:B\to B'$, and let $h:E\to E'$ be the context extension arrow. Thus $\pi'h=b\pi$. Compose the cartesian arrow $c:\pi^*Q\to Q$ with $Q\to Q'$. Its base is $b\pi=\pi'h$, so the cartesian arrow $c':\pi'^*Q'\to Q'$ uniquely factors it through a total arrow
\[
 \pi^*Q\longrightarrow\pi'^*Q'\quad\text{over }h.
\]
Identity and composition of these maps follow from that same uniqueness, not from faithfulness of the content extension. Naturality of the occurrence map preserves the actual addressed shape under $dQ\to dQ'$ while relabeling each position by $b$. The actual context fibre maps of $P\to P'$ therefore give the output context map position by position. Their coefficient/effect-map functoriality and the display square prove the full native body square; no equal-label address identification is made.

The input body $A\to dQ$ and output body $E\to C$ of a general representative then use context reindexing and content optransport. Their chosen factorization properties make the result functorial for the whole representative square. Pure-middle natural transformations use the original canonical Beck--Chevalley mate. Mate naturality for these total arrows and mate pasting supply its naturality and balance coherence along $J$. This verifies the extra presentation hypothesis for the actual O1 construction using the mechanisms in its native-cell, endpoint, and pseudobalance proof, rather than postulating new mixed comparisons. The four existing modifications are natural for these maps by their established coefficient and cartesian-grade comparisons.

This argument also applies to the alternative effect-fibre context input where the same native addressed crossing has been established. It is a closure result for those presentations, not for an arbitrary unpresented relative algebra.

\section{Change of diagram shape}
For a functor $r:I\to J$, precomposition gives pure, context and content functors from the $J$ construction to the $I$ construction. They preserve the extension and full native body square on the nose. Units, accumulation and crossing commute up to the original chosen cleavage comparisons, which are coherent pointwise. On coend representatives the map takes the middle diagram $Y$ to $Y r$ and takes a balance transformation to its restriction. Hence these maps preserve the presented relative multiplication; no assertion that all $I$-middle diagrams extend to $J$ is needed. Identity and composite shape changes have the evident coherent comparisons.

The terminal shape recovers the original presentation. A discrete shape remembers a family without comparison arrows; a star shape retains root-to-component native comparisons. Forgetting an arrow imposes fewer naturality constraints and need not have a reverse extension. Thus shape change is not an assertion that arbitrary missing dependence or compatibility data can be reconstructed.

\section{Dependence interpretation and exact limits}
For a star-shaped $J$ with arrows from a root to its components, a context diagram can retain an actual joint effect context and native marginal projections when the original context doctrine is an effect-fibre realization. Its pure and content diagrams retain the matching interface and body squares. The new crossing transforms this entire diagram coherently, rather than treating equal marginal laws as an equality of raw witnesses. Arbitrary shapes $J$ retain their specified relations; the theorem does not assert that every desired diagram exists.

The tree-overlap admission theorem can supply some such context diagrams in weak-pullback effect models. It is logically separate from the present closure theorem. Neither gives admission for arbitrary cyclic constraints. Conditional component behavior must still factor through its declared interface. Consequently arbitrary physical interference, arbitrary information hiding, scheduler availability, fairness and infinite effect laws are not supplied by diagram lifting.

This is not the cone multicategory with its opposite unary category: native arrows here keep their original direction. It is also not a new proof that every categorical law admits the specified representative lift. The hypothesis has to be verified on actual mixed middle/endpoint arrows; if only fixed-endpoint naturality is available, the proof stops at that step.
\section{Source comparison and attribution}
Functor-category lifting of fibred structure and componentwise coherence are standard categorical techniques. The project-specific task here is to verify the actual total representative crossing and reprove retained normalization over the diagram pure base; neither follows from an unsupported claim that diagram categories commute with pseudocoends. This is a closure theorem for supplied presentations, not a priority claim for a new distributive-law formalism.

The pinned internal source is \emph{General Independent Interaction Criterion: Augmented rig contexts and relative occurrence displayed content}, 1 November 2026, specifically its endpoint-square convention, retained comma normalization, construction of all native/pure crossing maps, and pure-middle Beck--Chevalley balance. The alternative input is the project's \emph{Independent Occurrence Interaction for Effect Fibres}. The pseudodistributive comparison format and composite machinery remain those already used there: Gambino and Lobbia, \emph{On the formal theory of pseudomonads and pseudodistributive laws}, TAC37(2),2021, and Marmolejo, \emph{Distributive laws for pseudomonads}, TAC5(5),1999. No new claim about the content of those references beyond the inherited framework is needed for the direct generator-level proof above.
\end{document}
